\chapter{Conclusions and Recommendations} \label{chap:Conclusion}
This study addresses the limitation of sequential solvers in terms of their computational costs and success rate especially when being applied to very complex combinatorial solution spaces. The results from a series of structured experiments showed that the multithreaded multi-colony ant optimization with ring and random communication topologies framework successfully outperforms the state-of-the-art solvers in terms of speed, reliability, and scalability. Although heterogenous ACS and MMAS colonies of the single-threaded multi-colony ACO solver help improve the diversity and convergence speed of the solver for smaller Sudoku puzzles. With its entropy-aware colonies, the CP-DCM-ACO solver was able to overcome the stagnation and exponential time increase suffered by the single colony ACO solver. The implementation of injecting the MMAS diversity to the ACS colony that is already stagnating and recommending paths to the MMAS colony when it converges too slowly successfully balances the exploration and exploitation of the ants. 

However, the solution space grows combinatorially as the grid size becomes larger. The collaboration between multiple colonies is still not enough to navigate a complex search space from high-order Sudoku puzzles like $25\times25$. By utilizing multiple threads running in parallel and placing this multi-colony ACO inside each thread, more regions of a very complex solution space are explored. But threads running in parallel without information sharing could lead to redundant visits to a certain region of the solution space. By allowing the threads to communicate in SudoSLVRR, the solver was able to prevent premature convergence by sending locally promising solutions from one thread to its neighbor. Also, the random communication topology allows the threads to receive promising solutions from the threads that explored the other region of the solution space thereby increasing the diversity. This shows that multi-level collaboration becomes necessary when the solution space expands exponentially. Overall, this research shows that the proposed approach effectively balances exploration and exploitation through multi-level collaboration, leading to a better performance across evaluated Sudoku instances. 

While this study establishes a new benchmark for solving high-order Sudoku, future study may explore other communication topologies other than ring and random communication topologies. SudoSLVRR executes the multi-colony ACO sequentially inside each thread and therefore accommodates further exploration in executing it in parallel. As GPU gains recognition when it comes to execution time, future study might explore implementing the multithreaded Sudoku solvers using GPU. Another promising area to explore for future work is the incorporation of negative learning in the ACO framework like in~\cite{Masukane, nurcahyadi2022negative, Haynes, james} to discourage exploration of dead-end solutions, thereby improving the solver’s performance.  Lastly, SudoSLVRR was applied to a Sudoku puzzle with standard constraints. Future researchers may explore applying the multithreaded solver to puzzles with improved constraints like 3D Sudoku and Diagonal Sudoku. 

